
\setkomafont{sectioning}{\rmfamily\bfseries} % Überschriften mit Serifen


% \ifenglish
\usepackage[ngerman,english]{babel}	    		% Überschriften werden von LaTeX in der korrekten Sprache erzeugt (english = Hauptsprache)
% \else
% \usepackage[english,ngerman]{babel}	    		% Überschriften werden von LaTeX in der korrekten Sprache erzeugt (ngerman = Hauptsprache)
% \fi

%\usepackage[ansinew]{inputenc}			% Umlaute
%\usepackage[utf8]{inputenc}
%\usepackage{subfig}
\usepackage[font=small,labelfont=bf]{subcaption}
\usepackage{placeins}
\usepackage{preview}
\newcommand{\kreis}[1]{\unitlength1ex\begin{picture}(2.5,2.5)\put(0.75,0.75){\circle{2.5}}\put(0.75,0.75){\makebox(0,0){#1}}\end{picture}}
\usepackage{diagbox}
\usepackage{colortbl}
\usepackage[T1]{fontenc}		        % Schriftart
\usepackage[utf8]{inputenc}	  		% Direkte Eingabe der Sonderzeichen über die Tastatur 
%\usepackage{lmodern}					% schönere Schrift
%\usepackage{ae}							% schöne Schrift
%\usepackage[final]{microtype}			% mikrotypografische Schriftanpassung
\usepackage{lipsum}
%\DisableLigatures{}

%\usepackage{fixltx2e}
\usepackage{graphicx}	
\usepackage{epstopdf}	
\usepackage{pgf}	       		% Ermöglicht die Einbindung von Bildern
\providecommand{\mathdefault}[1]{#1}	% benötigt für matplotlib pgf-Export
\usepackage{makeidx}		          	% Erzeugt Verzeichnisse
\usepackage{array}
\usepackage{float}		             	% Erzwingt Positionierung von Floatobjekten mit "`H"'
\usepackage{tabularx}		          	% Tabellen mit Spalten gleicher Breite
%\usepackage{tabu}						% sehr flexible Tabellen mit mehr Möglichkeiten
%\usepackage{dcolumn}		          	% Tabellen mit ausgerichteten Einträgen
\usepackage{amssymb,amsmath}	    	% Um Texte in mathematischen Gleichungen zu schreiben mit dem Befehl \text{}
%\usepackage[figuresright]{rotating}	% Damit können Bilder, Tabellen, usw. in landscape Format angezeigt werden
\usepackage{bibgerm}		           	% Ermöglicht das Einbinden deutschsprachiger Literaturverzeichnisse
%\usepackage{subfig}		            % Mehrere Bilder nebeneinander, unter gleicher Bildunterschrift und -nummer
\usepackage{booktabs}		          	% Hochwertige Tabellen
\usepackage{longtable}	        		% Tabellen über mehrere Seiten
\usepackage{url}			            % Weblinks
\usepackage{pdfpages}					% PDF einbinden
%\usepackage{eurosans}
\usepackage[bookmarks,pagebackref=false,pdftex,bookmarksopen=true,bookmarksnumbered]{hyperref}	% PDF-Bookmarks, Angabe im Literaturverzeichnis auf welcher Seite Einordnungsformel steht momentan auf false gesetzt
\usepackage{multirow}
%\usepackage[square]{natbib}         	% nützliches Paket, falls man die Gestaltung des Literaturverzeichnisse und der Referenzen selbst übernehmen möchte
\usepackage{wrapfig}
\usepackage{tocvsec2}
\usepackage{ifthen}
\usepackage{xspace}						%Leerzeichen nach Anführungszeichen
%\usepackage[decimalsymbol=comma, expproduct=cdot]{siunitx}				%Einheiten im Mathemodus
%\usepackage{microtype,textcomp}			% bei siunitx
\usepackage{bm}
\usepackage{pgfplots}					%Matlab Diagramme in Latex
\def\rb#1{\rotatebox{90}{$\xleftarrow{#1}$}}
\newcommand{\tikzmark}[1]{\tikz[overlay, remember picture] \coordinate (#1);}


%---------------------------------------
% Glossary
%---------------------------------------
\usepackage[acronym, toc, nonumberlist]{glossaries}
\makeglossaries


